\section{The exact higher K-theory quantity behind the strict Selmer length}
\label{sec:selmer-k-theory}

The finiteness in AC-G22 can be made into an exact order identity.
This identifies a concrete arithmetic quantity to estimate; it does
not replace an estimate by finiteness.  Throughout this section
$e\geq1$, and $F$ is a totally imaginary number field.  Put
\[
 A_{\ell,e}=\mathbb Q_\ell/\mathbb Z_\ell(-e),\qquad
 S_\ell(F,e)=\ker\left(H^1(G_{F,S_\ell},A_{\ell,e})
       \longrightarrow\bigoplus_{v\mid\ell}
                         H^1(F_v,A_{\ell,e})\right),
\]
where $S_\ell$ comprises the primes above $\ell$ and the complex
places.  Complex places contribute zero modified local cohomology.
This is the untwisted strict group used in AM-H06--AM-H08.
Write $B^\vee=\operatorname{Hom}_{\rm cont}(B,\mathbb Q_\ell/\mathbb Z_\ell)$
for Pontryagin duality on the relevant $\ell$-primary groups.

\begin{proposition}[AC-K10: strict order and the local correction]
\label{AC-K10}
There is an exact sequence of finite groups
\[
 0\longrightarrow S_\ell(F,e)^\vee
 \longrightarrow K_{2e}(\mathcal O_F)_{(\ell)}
 \longrightarrow\bigoplus_{v\mid\ell}
                     H^0(F_v,A_{\ell,e})^\vee
 \longrightarrow H^0(F,A_{\ell,e})^\vee
 \longrightarrow0.                                      \tag{K10.1}
\]
The middle identification is the standard \'etale Chern-character
identification; the last arrow is dual to restriction of invariants.
In particular
\[
 \#S_\ell(F,e)=
 \#K_{2e}(\mathcal O_F)_{(\ell)}\,
 \frac{\#H^0(F,A_{\ell,e})}
      {\prod_{v\mid\ell}\#H^0(F_v,A_{\ell,e})}.          \tag{K10.2}
\]
All these groups are finite, including when $\ell=2$.  For fixed
$F,e$, only finitely many $S_\ell(F,e)$ are nonzero, and
\[
 \sum_\ell \log\ell\,
       \operatorname{length}_{\mathbb Z_\ell}S_\ell(F,e)
 \leq\log\#K_{2e}(\mathcal O_F).                         \tag{K10.3}
\]
More generally, for every integer $t\geq2$, the left side summed
over $1\leq e<t$ is at most
$\sum_{1\leq e<t}\log\#K_{2e}(\mathcal O_F)$.
\end{proposition}
\begin{proof}
Take the compact representation $T=\mathbb Z_\ell(e+1)$.
Its Tate dual $T^\vee(1)$ is exactly $A_{\ell,e}$.
The compact form of Poitou--Tate duality identifies the kernel
of localization in $H^2(G_{F,S_\ell},T)$ with
$S_\ell(F,e)^\vee$.  The final part of the Poitou--Tate sequence is
\[
 H^2(G_{F,S_\ell},T)\longrightarrow
 \bigoplus_{v\mid\ell}H^2(F_v,T)\longrightarrow
 H^0(F,A_{\ell,e})^\vee\longrightarrow0.                 \tag{K10.4}
\]
Local Tate duality identifies $H^2(F_v,T)$ with
$H^0(F_v,A_{\ell,e})^\vee$, with the maps in (K10.4) dual
to the usual restrictions.  The exact statements used are
R.~Sharifi, \emph{Iwasawa Theory}, Appendix~A, Theorems
A.0.3, A.0.6 and A.0.8, with the compact-coefficient extensions
in Remarks A.0.4 and A.0.9.  Their use here includes continuous
compact cohomology, not an unjustified replacement of finite
coefficients in a finite-level sequence.  Because $F$ has no
real places, there is no real-place modification at $\ell=2$.

The integral \'etale comparison gives
\[
 H^2(G_{F,S_\ell},\mathbb Z_\ell(e+1))
 =H^2_{\rm et}(\mathcal O_F[1/\ell],\mathbb Z_\ell(e+1))
 \simeq K_{2e}(\mathcal O_F)_{(\ell)}.                  \tag{K10.5}
\]
Here the first equality is the usual arithmetic Galois-cohomology
description.  The precise K-theory result is C.~Weibel,
\emph{The K-book}, Chapter~VI, Theorem~8.2, with its totally
imaginary formulation in Theorem~8.4 (author's online chapter,
pp.~47--48).  It applies for $2e\geq2$ and includes $\ell=2$
under the present hypothesis.  These are cited known theorems;
their proofs are not claimed as new results of this run.
The even group $K_{2e}(\mathcal O_F)$ is finite.  Substitution
in (K10.4), together with the kernel identification, gives
(K10.1).

For explicit finiteness of the invariant terms, the cyclotomic
image of an absolute number-field or $\ell$-adic local Galois
group is an open subgroup of $\mathbb Z_\ell^\times$.
It contains an element $u$ of infinite order.  Since $e>0$,
$u^{-e}-1\ne0$, and the invariant subgroup of
$\mathbb Q_\ell/\mathbb Z_\ell(-e)$ is killed by that nonzero
scalar.  It is consequently finite.  Taking alternating orders
in (K10.1) proves (K10.2).

Restriction injects global invariants into the invariants at
each chosen place: both are subgroups of the same coefficient
module, and the global Galois group contains its decomposition
group.  In particular the quotient of invariant orders in
(K10.2) is at most one.  Thus the order of $S_\ell(F,e)$ divides
the order of the indicated $\ell$-primary K-group; the injection
in (K10.1) gives the same conclusion.  All but finitely many
such K-groups are zero.  Taking logarithms and summing proves
(K10.3) and its finite-weight version.
\end{proof}

\begin{proposition}[AC-K11: the exact local subtraction for the cubic field]
\label{AC-K11}
Let $K=\mathbb Q(\omega)$ and $E=K(\sqrt[3]{269})$.  Define
\[
 \kappa_e=v_3\bigl(\#K_{2e}(\mathcal O_E)\bigr),\qquad
 k_e=1+v_3(e),\qquad
 a_E(e)=\operatorname{length}_{\mathbb Z_3}S_3(E,e).
\]
Then
\[
                a_E(e)=\kappa_e-2k_e.                  \tag{K11.1}
\]
The corresponding strict group over $K$ is zero.  Hence the
finite cyclic exponent $m_e$ of AC-G22 satisfies
\[
 \kappa_e-2k_e-4\ \leq\ m_e\ \leq\
                       \kappa_e-2k_e+9.               \tag{K11.2}
\]
The lower bound may of course be strengthened to $m_e\geq1$.
These identities and inequalities use no RH hypothesis.
They reduce the positive-weight coefficient-prime-$3$ estimate
to the growth of a definite sequence of even K-groups, up to
an explicit logarithmic local term and a bounded integral
character-comparison error.  They do not supply a growth bound
for $\kappa_e$ or for the fields appearing in the source's
growing-host family.
\end{proposition}
\begin{proof}
The Hensel argument in AC-G22 proves that $269$ is a cube in
$\mathbb Q_3$.  The three roots of $X^3-269$ are then in
$K_3=\mathbb Q_3(\omega)$, so the unique prime of $K$ over $3$
splits into three primes of $E$, each with completion $K_3$.
At the prime $(269)$ of $K$ the cubic extension $E/K$ is
totally ramified, while the cyclotomic $3$-power tower is
unramified there.  Since $[E:K]=3$, its intersection with
$\mathbb Q(\mu_{3^\infty})$ is therefore exactly $K$.
It follows that the $3$-adic cyclotomic images for $G_E$ and
for all three completions are the same group
$1+3\mathbb Z_3$.

For that group the minimum of
$v_3(u^{-e}-1)$ is $k_e=1+v_3(e)$, by the binomial
lifting-of-exponent calculation; $u=4$ attains it.
Thus each of the four invariant groups in (K10.2), one global
and three local, has order $3^{k_e}$.  Formula (K10.2)
therefore gives
$3^{a_E(e)}=3^{\kappa_e}3^{k_e}/3^{3k_e}$,
which is (K11.1).

For $K$, one can use either the direct residual strict-kernel
calculation accompanying AC-G22 or the following known
K-theory computation.  The prime $3$ is regular because
$\mathbb Z[\omega]$ has class number one.  Weibel, Chapter~VI,
Example~8.3.2, gives
$K_{2e}(\mathbb Z[\omega])_{(3)}=0$ for all $e\geq1$.
The injection in (K10.1) proves $S_3(K,e)=0$.

The existing integral comparison AM-H08, (H08.2)--(H08.3),
for $r=3,p=269$, reads
\[
       -9\leq a_E(e)-a_K(e)-m_e\leq4.                 \tag{K11.3}
\]
All its finiteness hypotheses now hold by AC-K10 and AC-G22.
Its $S$ contains the primes above $3$ and $269$; its strict
inertia conditions identify the untwisted groups with
$S_3(E,e)$ and $S_3(K,e)$, as proved in AM-H06.  Thus the
comparison concerns the very groups used here, rather than
groups with a ramified local condition omitted.  Substitution
of $a_K(e)=0$ and (K11.1), followed by rearrangement, gives
(K11.2).
\end{proof}

\begin{proposition}[AC-K12: the single cyclic K-module at three]
\label{AC-K12}
For $E=\mathbb Q(\omega,\sqrt[3]{269})$,
$\operatorname{Cl}(E)_{(3)}=0$.
Let $\Gamma=\operatorname{Gal}(E/K)=\langle\sigma\rangle$.
Via $\sigma\mapsto\omega$, the action on
$H_e=K_{2e}(\mathcal O_E)_{(3)}$ factors through
$\mathcal O=\mathbb Z_3[\omega]$, and
\[
 H_e\simeq\mathcal O/\lambda^{\kappa_e}\mathcal O,
 \qquad
 S_3(E,e)^\vee\simeq
       \mathcal O/\lambda^{\kappa_e-2k_e}\mathcal O.     \tag{K12.1}
\]
The second isomorphism uses the induced $\Gamma$ action under
Poitou--Tate duality; length zero means the zero module.
In particular $\kappa_e\geq2k_e$.
This identifies the module structure, but gives no upper bound
on its exponent.
\end{proposition}
\begin{proof}
The only finite prime ramified in $E/K$ is $(269)$, with
ramification index three, by AC-K04; the primes over three
split by AC-G22.  There are no real places.
Chevalley's ambiguous class number formula gives
\[
 \#\operatorname{Cl}(E)^\Gamma
 =\frac{h_K\prod_v e_v}
 {[E:K]\,[\mathcal O_K^\times:
                 \mathcal O_K^\times\cap N_{E/K}E^\times]}
 =\frac1{[\mathcal O_K^\times:
                 \mathcal O_K^\times\cap N_{E/K}E^\times]}.
                                                               \tag{K12.2}
\]
The formula used is C.-F.~Yu, \emph{Ambiguous class number
formulas}, MPIM Preprint 2014(68), Theorem~1.1, equation~(1.2),
printed p.~1, with the ordinary-class interpretation on p.~2.
Its left side is a positive integer.  Thus the unit index and
the ambiguous class group order are both one.  If the finite
$3$-group $\operatorname{Cl}(E)_{(3)}$ were nonzero, its
nonzero $3$-torsion vector space would have a nonzero fixed
vector: on a vector space in characteristic three,
$(\sigma-1)^3=\sigma^3-1=0$ has nontrivial kernel.
This contradicts (K12.2), proving the first assertion.

We next justify the action on the K-group, including the
ramified prime.  Transfer along the finite free extension
$\mathcal O_K\subset\mathcal O_E$ sends $H_e$ to
$K_{2e}(\mathcal O_K)_{(3)}=0$, as used in AC-K11.
Transfer followed by restriction on the fraction fields is
$1+\sigma+\sigma^2$.  The localization sequence injects
$K_{2e}(\mathcal O_E)$ into $K_{2e}(E)$, since the groups
$K_{2e}$ of its finite residue fields vanish for $e>0$.
Consequently the same norm operator is zero on $H_e$.
This proves that its $\mathbb Z_3[\Gamma]$ action factors
through $\mathbb Z_3[\Gamma]/(1+\sigma+\sigma^2)
\simeq\mathcal O$.  This transfer argument does not assume
that $\operatorname{Spec}\mathcal O_E\to\operatorname{Spec}\mathcal O_K$
is \'etale at $269$.

The mod-three comparison of Weibel, Chapter~VI, Corollary~8.3
and Example~8.3.1, gives
\[
 \dim_{\mathbb F_3}(H_e/3H_e)
 =\dim_{\mathbb F_3}\bigl(\operatorname{Cl}
                         (\mathcal O_E[1/3])/3\bigr)
            +\#\{v\mid3\}-1=2.                       \tag{K12.3}
\]
The localized class group has no three-part, being a quotient
of $\operatorname{Cl}(E)$.  There are exactly three places
over three.

The local quotient in (K10.1) can be identified equivariantly.
Choose a lift of $\sigma$ to $G_K$ with trivial $3$-adic
cyclotomic character; it exists because
$E\cap\mathbb Q(\mu_{3^\infty})=K$.
It cyclically permutes the three places above three and acts
trivially on their common invariant coefficients.  The last
map in (K10.1) is consequently the sum
\[
 (\mathbb Z/3^{k_e})^3\longrightarrow\mathbb Z/3^{k_e}.
\]
Its kernel is the augmentation module for $\Gamma$ over
$\mathbb Z/3^{k_e}$.  The map $a\mapsto(\sigma-1)a$ induces
an isomorphism from $\mathcal O=\mathbb Z_3[\Gamma]/(1+\sigma+\sigma^2)$
onto the integral augmentation ideal; reduction modulo $3^{k_e}$ therefore
identifies this kernel with $\mathcal O/3^{k_e}\mathcal O$.
Thus (K10.1) gives a surjection
\[
              H_e\longrightarrow
                 \mathcal O/\lambda^{2k_e}\mathcal O.  \tag{K12.4}
\]

Write the finite $\mathcal O$-module as
$H_e=\bigoplus_i\mathcal O/\lambda^{n_i}$ with $n_i>0$.
Equation (K12.3) says $\sum_i\min(2,n_i)=2$.
The quotient (K12.4), with $2k_e\geq2$, forces at least
one $n_i\geq2k_e$.  Its summand already contributes two
to that sum, so there are no other summands.
Its length is $\kappa_e$, since the residue field has order
three.  The kernel of the surjection (K12.4) from this
cyclic module is cyclic of length $\kappa_e-2k_e$.
By (K10.1) it is precisely $S_3(E,e)^\vee$, proving (K12.1).
\end{proof}

\paragraph{Exact remaining size statement.}
For any totally imaginary host field $F$ and integer $t$, AC-K10
supplies the concrete sufficient estimate
\[
 \sum_{1\leq e<t}\log\#K_{2e}(\mathcal O_F)\leq B
 \quad\Longrightarrow\quad
 \sum_\ell\log\ell\sum_{1\leq e<t}
       \operatorname{length}_{\mathbb Z_\ell}S_\ell(F,e)
       \leq B.                                        \tag{K11.4}
\]
The quantifiers over the fields, weights and allowed bound in
the source predicate \texttt{VVle} must still be satisfied.
An unspecified finite bound for each fixed $F,t$ is not a
uniform bound with the source's prescribed $B$.
Nor is an estimate over the fixed field $E$ automatically an
estimate over the varying CM hosts in \texttt{SupplyChi}.
This is the exact limitation of the K-theory route at this stage.

\paragraph{Algebraic-geometric dictionary.}
The Kummer character is the lisse rank-one coefficient system
of the finite \'etale torsor over
$\operatorname{Spec}\mathcal O_K[1/(3\cdot269)]$ (SGA~1,
XI~\S6, Propositions~1 and~4 and Corollary~5).  AC-K10 places
its untwisted restriction in localization and duality on the
arithmetic scheme $\operatorname{Spec}\mathcal O_E[1/\ell]$;
AC-K11 gives the numerical correction for its three local
fibers over $3$.  The passage from algebraic K-theory to these
cohomology groups is the modern \'etale comparison just cited,
not a theorem attributed to a Grothendieck archive scan.
The exterior-power and coefficient-tensor language belongs to
SGA~6, V~2.2(b), whereas the order estimate in (K11.4) is an
additional arithmetic requirement.  No previously inaccessible
Grothendieck scan is represented as reread by this argument.
